% This file is generated from the corresponding Obsidian note.
% Source: Teoricos/GS - Teo20.md
\chapter{Álgebra multilineal y formas alternantes}
\label{ch:teo-20}
\subsection{Un poco de algebra multilineal}
\begin{bookdefinition}{}
Sean $V_1,V_2,\ldots,V_k$ y $W$ espacios vectoriales. Una funcion

\[
F:V_1\times\cdots\times V_k\to W
\]
se dice multilineal si esta es lineal en cada variable, es decir, $\lambda\in\mathbb R$,

\[
F(v_1,\ldots,\lambda v_i+\tilde v_i,\ldots,v_k) = \lambda F(v_1,\ldots,v_i,\ldots,v_k)+F(v_1,\ldots,\tilde v_i,\ldots,v_k),
\]
$i=1,\ldots,k$.

\end{bookdefinition}
\phantomsection\label{blk:2f1d75}
\begin{bookdefinition}{Conjunto de funciones multilineales}
Se suele denotar por $L(V_1,\ldots,V_k;W)$ el conjunto de todas las funciones multilineales de $V_1\times\cdots\times V_k$ en $W$, el cual resulta ser un espacio vectorial real bajo la suma y producto escalar natural. Si $\lambda\in\mathbb R$:

\[
(F+G)(v_1,\ldots,v_k):=F(v_1,\ldots,v_k)+G(v_1,\ldots,v_k),
\]
\[
(\lambda F)(v_1,\ldots,v_k):=\lambda F(v_1,\ldots,v_k).
\]
\end{bookdefinition}
\begin{bookexample}{}
\[
\det:\underbrace{\mathbb R^n\times\cdots\times\mathbb R^n}_{n\text{-veces}}\to\mathbb R
\]
\[
\det(v_1,\ldots,v_n)=\det(v_1|\cdots|v_n).
\]
\end{bookexample}
\begin{bookexample}{}
\[
\langle\ ,\ \rangle:\mathbb R^n\times\mathbb R^n\to\mathbb R
\]
\[
((x_1,\ldots,x_n),(y_1,\ldots,y_n))\mapsto \langle v_1,v_2\rangle=\sum_{k=1}^n x_k y_k.
\]
Producto interno usual de $\mathbb R^n$.

\end{bookexample}
\begin{bookexample}{}
\[
\times:\mathbb R^3\times\mathbb R^3\to\mathbb R^3
\]
\[
((x_1,x_2,x_3),(y_1,y_2,y_3)) \mapsto \det\begin{pmatrix} e_1&e_2&e_3\\ x_1&x_2&x_3\\ y_1&y_2&y_3 \end{pmatrix},
\]
el producto cruz en $\mathbb R^3$.

\end{bookexample}
\begin{bookdefinition}{Producto tensorial}
Sean $V_1,\ldots,V_k,W_1,\ldots,W_l$ espacios vectoriales reales y sean

\[
F\in L(V_1,\ldots,V_k;\mathbb R),\qquad G\in L(W_1,\ldots,W_l;\mathbb R),
\]
y considerar la funcion

\[
F\otimes G:V_1\times\cdots\times V_k\times W_1\times\cdots\times W_l\to\mathbb R
\]
\[
(v_1,\ldots,v_k,w_1,\ldots,w_l)\mapsto F(v_1,\ldots,v_k)G(w_1,\ldots,w_l).
\]
Es facil ver que $F\otimes G$ es multilineal, asi que

\[
F\otimes G\in L(V_1,\ldots,V_k,W_1,\ldots,W_l;\mathbb R)
\]
y se llama el producto tensorial de $F$ y $G$.

\end{bookdefinition}
\phantomsection\label{blk:db8d92}
\begin{bookexercise}{}
Probar que el producto tensorial es bilineal y asociativo:

\[
(F+\lambda H)\otimes G=F\otimes G+\lambda(H\otimes G),
\]
\[
F\otimes(G+\beta K)=F\otimes G+\beta(F\otimes K),
\]
\[
(F\otimes G)\otimes K=F\otimes(G\otimes K).
\]
\end{bookexercise}
\begin{bookproposition}{Ejercicio}
Sean $V_1,\ldots,V_k$ espacios vectoriales reales de dimension finita $n_1,\ldots,n_k$ respectivamente, y sea $\{e^j_1,\ldots,e^j_{n_j}\}$ una base de $V_j$, $j=1,\ldots,k$, y sean $\{\varepsilon^j_1,\ldots,\varepsilon^j_{n_j}\}$ la respectiva base dual de $V_j^*$.

Entonces

\[
B=\left\{\varepsilon^1_{i_1}\otimes\cdots\otimes\varepsilon^k_{i_k}:1\le i_j\le n_j,\ 1\le j\le k\right\}
\]
es una base para $L(V_1,\ldots,V_k;\mathbb R)$, el cual es un espacio vectorial real de dimension $n_1n_2\cdots n_k$.

\end{bookproposition}
\begin{bookproof}{hecha por chat}
\begin{enumerate}
\item Para cada multiíndice $(i_1,\ldots,i_k)$, definimos
\end{enumerate}
\[
E_{i_1,\ldots,i_k}:=\varepsilon^1_{i_1}\otimes\cdots\otimes\varepsilon^k_{i_k}.
\]
\begin{enumerate}[start=2]
\item Por definición del producto tensorial, para $v_j\in V_j$, se tiene
\end{enumerate}
\[
E_{i_1,\ldots,i_k}(v_1,\ldots,v_k)=\varepsilon^1_{i_1}(v_1)\cdots\varepsilon^k_{i_k}(v_k).
\]
\begin{enumerate}[start=3]
\item En particular, si evaluamos en vectores de las bases, obtenemos
\end{enumerate}
\[
E_{i_1,\ldots,i_k}(e^1_{a_1},\ldots,e^k_{a_k})=\varepsilon^1_{i_1}(e^1_{a_1})\cdots\varepsilon^k_{i_k}(e^k_{a_k}).
\]
\begin{enumerate}[start=4]
\item Como $\{\varepsilon^j_1,\ldots,\varepsilon^j_{n_j}\}$ es la base dual de $\{e^j_1,\ldots,e^j_{n_j}\}$, vale
\end{enumerate}
\[
\varepsilon^j_{i_j}(e^j_{a_j})=\delta_{i_ja_j}.
\]
\begin{enumerate}[start=5]
\item Luego
\end{enumerate}
\[
E_{i_1,\ldots,i_k}(e^1_{a_1},\ldots,e^k_{a_k})=\delta_{i_1a_1}\cdots\delta_{i_ka_k}.
\]
\begin{enumerate}[start=6]
\item Por lo tanto, este valor es $1$ si $(i_1,\ldots,i_k)=(a_1,\ldots,a_k)$, y es $0$ en caso contrario.
\item Probemos primero que $B$ es linealmente independiente. Supongamos que
\end{enumerate}
\[
\sum_{i_1=1}^{n_1}\cdots\sum_{i_k=1}^{n_k}c_{i_1,\ldots,i_k}E_{i_1,\ldots,i_k}=0.
\]
\begin{enumerate}[start=8]
\item Evaluamos ambos lados en $(e^1_{a_1},\ldots,e^k_{a_k})$. Entonces
\end{enumerate}
\[
0=\sum_{i_1=1}^{n_1}\cdots\sum_{i_k=1}^{n_k}c_{i_1,\ldots,i_k}E_{i_1,\ldots,i_k}(e^1_{a_1},\ldots,e^k_{a_k}).
\]
\begin{enumerate}[start=9]
\item Usando el cálculo anterior,
\end{enumerate}
\[
0=\sum_{i_1=1}^{n_1}\cdots\sum_{i_k=1}^{n_k}c_{i_1,\ldots,i_k}\delta_{i_1a_1}\cdots\delta_{i_ka_k}.
\]
\begin{enumerate}[start=10]
\item Todos los términos se anulan salvo el término con $(i_1,\ldots,i_k)=(a_1,\ldots,a_k)$, por lo tanto
\end{enumerate}
\[
c_{a_1,\ldots,a_k}=0.
\]
\begin{enumerate}[start=11]
\item Como $(a_1,\ldots,a_k)$ era arbitrario, todos los coeficientes son cero. Luego $B$ es linealmente independiente.
\item Probemos ahora que $B$ genera $L(V_1,\ldots,V_k;\mathbb R)$. Sea $F\in L(V_1,\ldots,V_k;\mathbb R)$. Queremos probar que
\end{enumerate}
\[
F=\sum_{i_1=1}^{n_1}\cdots\sum_{i_k=1}^{n_k}F(e^1_{i_1},\ldots,e^k_{i_k})E_{i_1,\ldots,i_k}.
\]
\begin{enumerate}[start=13]
\item Sean $v_j\in V_j$ arbitrarios. Como $\{e^j_1,\ldots,e^j_{n_j}\}$ es base de $V_j$, existen escalares $a^j_{i_j}$ tales que
\end{enumerate}
\[
v_j=\sum_{i_j=1}^{n_j}a^j_{i_j}e^j_{i_j}.
\]
\begin{enumerate}[start=14]
\item Como $\{\varepsilon^j_1,\ldots,\varepsilon^j_{n_j}\}$ es la base dual, los coeficientes están dados por
\end{enumerate}
\[
a^j_{i_j}=\varepsilon^j_{i_j}(v_j).
\]
\begin{enumerate}[start=15]
\item Entonces, por multilinealidad de $F$,
\end{enumerate}
\[
F(v_1,\ldots,v_k)=F\left(\sum_{i_1=1}^{n_1}a^1_{i_1}e^1_{i_1},\ldots,\sum_{i_k=1}^{n_k}a^k_{i_k}e^k_{i_k}\right).
\]
\begin{enumerate}[start=16]
\item Expandiendo en cada variable, obtenemos
\end{enumerate}
\[
F(v_1,\ldots,v_k)=\sum_{i_1=1}^{n_1}\cdots\sum_{i_k=1}^{n_k}a^1_{i_1}\cdots a^k_{i_k}F(e^1_{i_1},\ldots,e^k_{i_k}).
\]
\begin{enumerate}[start=17]
\item Sustituyendo $a^j_{i_j}=\varepsilon^j_{i_j}(v_j)$, queda
\end{enumerate}
\[
F(v_1,\ldots,v_k)=\sum_{i_1=1}^{n_1}\cdots\sum_{i_k=1}^{n_k}\varepsilon^1_{i_1}(v_1)\cdots\varepsilon^k_{i_k}(v_k)F(e^1_{i_1},\ldots,e^k_{i_k}).
\]
\begin{enumerate}[start=18]
\item Por definición de $E_{i_1,\ldots,i_k}$, esto es
\end{enumerate}
\[
F(v_1,\ldots,v_k)=\sum_{i_1=1}^{n_1}\cdots\sum_{i_k=1}^{n_k}F(e^1_{i_1},\ldots,e^k_{i_k})E_{i_1,\ldots,i_k}(v_1,\ldots,v_k).
\]
\begin{enumerate}[start=19]
\item Como esto vale para todo $(v_1,\ldots,v_k)$, concluimos que
\end{enumerate}
\[
F=\sum_{i_1=1}^{n_1}\cdots\sum_{i_k=1}^{n_k}F(e^1_{i_1},\ldots,e^k_{i_k})E_{i_1,\ldots,i_k}.
\]
\begin{enumerate}[start=20]
\item Luego $B$ genera $L(V_1,\ldots,V_k;\mathbb R)$.
\item Como $B$ es linealmente independiente y genera $L(V_1,\ldots,V_k;\mathbb R)$, $B$ es una base.
\item Finalmente, para elegir un elemento de $B$, debemos elegir un índice $i_j$ entre $1$ y $n_j$ para cada $j=1,\ldots,k$. Por lo tanto $B$ tiene $n_1n_2\cdots n_k$ elementos, y entonces
\end{enumerate}
\[
\dim L(V_1,\ldots,V_k;\mathbb R)=n_1n_2\cdots n_k.
\]
\bookqed
\end{bookproof}
\begin{bookdefinition}{}
Sean $V_1,\ldots,V_k$ espacios vectoriales y $V_1^*,\ldots,V_k^*$ sus respectivos espacios duales. Se denota por

\[
V_1^*\otimes\cdots\otimes V_k^*
\]
al subespacio de $L(V_1,\ldots,V_k;\mathbb R)$ generado por

\[
\{\alpha_1\otimes\cdots\otimes\alpha_k:\alpha_i\in V_i^*\}.
\]
Por la prop. anterior:

\[
L(V_1,\ldots,V_k;\mathbb R)=V_1^*\otimes\cdots\otimes V_k^*.
\]
\end{bookdefinition}
\begin{bookproposition}{Propiedad caracteristica del producto tensorial}
Sean $V_1,\ldots,V_k$ espacios vectoriales reales de dimension finita y sea $W$ un espacio vectorial real (no necesariamente de dimension finita). Sea

\[
A:V_1^*\times\cdots\times V_k^*\to W
\]
en $L(V_1^*,\ldots,V_k^*;W)$.

Entonces existe una unica transformacion lineal

\[
\hat A:V_1^*\otimes\cdots\otimes V_k^*\to W
\]
tal que $A=\hat A\circ\pi$, con

\[
\pi:V_1^*\times\cdots\times V_k^*\to V_1^*\otimes\cdots\otimes V_k^*,
\]
\[
\pi(\alpha_1,\ldots,\alpha_k)=\alpha_1\otimes\cdots\otimes\alpha_k.
\]
\end{bookproposition}
\begin{bookproof}{}
\begin{enumerate}
\item Fijar base $\{\varepsilon^i_1,\ldots,\varepsilon^i_{n_i}\}$ base dual de $V_i^*$.
\item La multilinealidad de $A$ nos dice que $A$ va estar determinada por lo que haga en la $k$-upla
\end{enumerate}
\[
(\varepsilon^1_{i_1},\ldots,\varepsilon^k_{i_k}).
\]
\begin{enumerate}[start=3]
\item Llamemos
\end{enumerate}
\[
w_{i_1,\ldots,i_k}=A(\varepsilon^1_{i_1},\ldots,\varepsilon^k_{i_k}).
\]
\begin{enumerate}[start=4]
\item Por la proposicion anterior, $\varepsilon^1_{i_1}\otimes\cdots\otimes\varepsilon^k_{i_k}$ es base de $V_1^*\otimes\cdots\otimes V_k^*$, entonces definimos $\hat A$ sobre los elementos de la base como
\end{enumerate}
\[
\hat A(\varepsilon^1_{i_1}\otimes\cdots\otimes\varepsilon^k_{i_k})=w_{i_1,\ldots,i_k}.
\]
\begin{enumerate}[start=5]
\item Y $\hat A$ es la unica que cumple esta condicion, pues cualquier otra, digamos $\tilde A$, al cumplir
\end{enumerate}
\[
\tilde A\circ\pi=A,
\]
va a implicar que

\[
\tilde A\circ\pi(\varepsilon^1_{i_1},\ldots,\varepsilon^k_{i_k})=w_{i_1,\ldots,i_k},
\]
\begin{enumerate}[start=6]
\item Es decir
\end{enumerate}
\[
\tilde A(\varepsilon^1_{i_1}\otimes\cdots\otimes\varepsilon^k_{i_k})=w_{i_1,\ldots,i_k}.
\]
\bookqed
\end{bookproof}
\section{Tensores covariantes}
\begin{bookdefinition}{$k$-tensores covariante}
Sea $V^*$ un espacio vectorial real de dimension finita y sea $k\in\mathbb N$. Un $k$-tensor covariante sobre $V$ es un elemento de

\[
\underbrace{V^*\otimes\cdots\otimes V^*}_{k\text{-veces}},
\]
lo cual es lo mismo a tener una funcion multilineal real valuada

\[
\alpha:\underbrace{V\times\cdots\times V}_{k\text{-veces}}\to\mathbb R.
\]
El numero $k$ se llama rango del tensor. Por convencion, un $0$-tensor es un numero real. Se denota al conjunto del los $k$ tensores por $T^k(V^*):=V^*\otimes\cdots\otimes V^*$, Entonces

\[
T^1(V^*)=V^*,\qquad T^0(V^*)=\mathbb R.
\]
\end{bookdefinition}
\begin{bookdefinition}{Tensores alternantes}
Sea $\alpha:V\times\cdots\times V\to\mathbb R$ un $k$-tensor covariante. Se dice que $\alpha$ es alternante, si para todo $\sigma\in S_k$, una permutacion de $k$ elementos, se cumple

\[
\alpha(v_{\sigma(1)},\ldots,v_{\sigma(k)})=\operatorname{sg}(\sigma)\alpha(v_1,\ldots,v_k)
\]
para todo $v_1,\ldots,v_k\in V$. Se denota por $\Lambda^k(V^*)$ el subconjunto de $T^k(V^*)$ formado solo por los que son alternantes, el cual es naturalmente un subespacio vectorial.

\end{bookdefinition}
\begin{bookexercise}{}
Mostrar que $\alpha\in\Lambda^k(V^*)$ sii

\[
\alpha(v_1,\ldots,v_i,\ldots,v_j,\ldots,v_k)=-\alpha(v_1,\ldots,v_j,\ldots,v_i,\ldots,v_k)
\]
para todo $i,j$ con $1\le i<j\le k$.

\end{bookexercise}
\begin{bookproof}{}
\begin{enumerate}
\item Supongamos primero que $\alpha\in\Lambda^k(V^*)$.
\item Por definición, para toda permutación $\sigma\in S_k$ se cumple
\end{enumerate}
\[
\alpha(v_{\sigma(1)},\dots,v_{\sigma(k)})=\operatorname{sg}(\sigma)\alpha(v_1,\dots,v_k).
\]
\begin{enumerate}[start=3]
\item Fijemos $i<j$ y sea $\tau\in S_k$ la transposición que intercambia $i$ y $j$, dejando fijos los demás índices.
\item Entonces $\operatorname{sg}(\tau)=-1$.
\item Aplicando la definición de tensor alternante a la transposición $\tau$, obtenemos
\end{enumerate}
\[
\alpha(v_{\tau(1)},\dots,v_{\tau(k)})=-\alpha(v_1,\dots,v_k).
\]
\begin{enumerate}[start=6]
\item Pero $(v_{\tau(1)},\dots,v_{\tau(k)})$ es la lista $(v_1,\dots,v_j,\dots,v_i,\dots,v_k)$, porque $\tau$ intercambia solamente las posiciones $i$ y $j$.
\item Por lo tanto,
\end{enumerate}
\[
\alpha(v_1,\dots,v_j,\dots,v_i,\dots,v_k)=-\alpha(v_1,\dots,v_i,\dots,v_j,\dots,v_k).
\]
\begin{enumerate}[start=8]
\item Equivalentemente,
\end{enumerate}
\[
\alpha(v_1,\dots,v_i,\dots,v_j,\dots,v_k)=-\alpha(v_1,\dots,v_j,\dots,v_i,\dots,v_k).
\]
\begin{enumerate}[start=9]
\item Recíprocamente, supongamos que para todo $i<j$ se cumple
\end{enumerate}
\[
\alpha(v_1,\dots,v_i,\dots,v_j,\dots,v_k)=-\alpha(v_1,\dots,v_j,\dots,v_i,\dots,v_k).
\]
\begin{enumerate}[start=10]
\item Queremos probar que $\alpha\in\Lambda^k(V^*)$. Es decir, queremos probar que para toda permutación $\sigma\in S_k$ se cumple
\end{enumerate}
\[
\alpha(v_{\sigma(1)},\dots,v_{\sigma(k)})=\operatorname{sg}(\sigma)\alpha(v_1,\dots,v_k).
\]
\begin{enumerate}[start=11]
\item Sea entonces $\sigma\in S_k$. Sabemos que toda permutación puede escribirse como producto de transposiciones, es decir, existen transposiciones $\tau_1,\dots,\tau_m$ tales que
\end{enumerate}
\[
\sigma=\tau_m\circ\cdots\circ\tau_1.
\]
\begin{enumerate}[start=12]
\item Cada transposición intercambia dos entradas. Por hipótesis, cada vez que intercambiamos dos entradas, el valor de $\alpha$ cambia de signo.
\item Aplicando sucesivamente esta propiedad a las transposiciones $\tau_1,\dots,\tau_m$, obtenemos
\end{enumerate}
\[
\alpha(v_{\sigma(1)},\dots,v_{\sigma(k)})=(-1)^m\alpha(v_1,\dots,v_k).
\]
\begin{enumerate}[start=14]
\item Por definición del signo de una permutación escrita como producto de transposiciones, se tiene
\end{enumerate}
\[
\operatorname{sg}(\sigma)=(-1)^m.
\]
\begin{enumerate}[start=15]
\item Luego
\end{enumerate}
\[
\alpha(v_{\sigma(1)},\dots,v_{\sigma(k)})=\operatorname{sg}(\sigma)\alpha(v_1,\dots,v_k).
\]
\begin{enumerate}[start=16]
\item Como esto vale para toda $\sigma\in S_k$, se sigue que $\alpha$ es alternante.
\item Por lo tanto, $\alpha\in\Lambda^k(V^*)$.
\end{enumerate}
\bookqed
\end{bookproof}
\begin{bookexample}{Algunos $k$-tensores alternantes}
\begin{itemize}
\item Cualquier $\alpha\in V^*=T^1(V^*)$ es alternante.
\item Tambien
\end{itemize}
\[
\det:\underbrace{\mathbb R^n\times\cdots\times\mathbb R^n}_{n\text{-veces}}\to\mathbb R
\]
esta en $\Lambda^n((\mathbb R^n)^*)$.

\end{bookexample}
\begin{bookdefinition}{Alternador}
Dado $\theta\in T^k(V^*)$, se define la parte antisimetrica de $\theta$, denotada por $\operatorname{Alt}(\theta)$, como el $k$-tensor

\[
\operatorname{Alt}(\theta)(v_1,\ldots,v_k)=\frac1{k!}\sum_{\sigma\in S_k}(\operatorname{sign}\sigma)\theta(v_{\sigma(1)},\ldots,v_{\sigma(k)}).
\]
\end{bookdefinition}
\begin{bookproposition}{Ejercicio}
\begin{itemize}
\item Para todo $\theta\in T^k(V^*)$, se tiene $\operatorname{Alt}(\theta)\in\Lambda^k(V^*)$, lo cual define una transformacion lineal
\end{itemize}
\[
\operatorname{Alt}:T^k(V^*)\to\Lambda^k(V^*).
\]
\begin{itemize}
\item Si $\theta\in\Lambda^k(V^*)$, entonces $\operatorname{Alt}(\theta)=\theta$.
\end{itemize}
\end{bookproposition}
\begin{bookproof}{}
\begin{itemize}
\item $\operatorname{Alt}(\theta)\in\Lambda^k(V^*)$
\end{itemize}
\begin{enumerate}
\item Recordemos que, para $\theta\in T^k(V^*)$, se define
\end{enumerate}
\[
\operatorname{Alt}(\theta)(v_1,\dots,v_k):=\frac{1}{k!}\sum_{\sigma\in S_k}\operatorname{sg}(\sigma)\theta(v_{\sigma(1)},\dots,v_{\sigma(k)}).
\]
\begin{enumerate}[start=2]
\item Primero probamos que $\operatorname{Alt}(\theta)$ es multilineal. Para cada $\sigma\in S_k$, la aplicación
\end{enumerate}
\[
(v_1,\dots,v_k)\mapsto \theta(v_{\sigma(1)},\dots,v_{\sigma(k)})
\]
es multilineal, porque $\theta$ es multilineal.\\
3. Como $\operatorname{Alt}(\theta)$ es una combinación lineal finita de aplicaciones multilineales, se sigue que $\operatorname{Alt}(\theta)$ es multilineal.\\
4. Ahora probamos que $\operatorname{Alt}(\theta)$ es alternante. Sea $\tau\in S_k$. Queremos probar que

\[
\operatorname{Alt}(\theta)(v_{\tau(1)},\dots,v_{\tau(k)})=\operatorname{sg}(\tau)\operatorname{Alt}(\theta)(v_1,\dots,v_k).
\]
\begin{enumerate}[start=5]
\item Por definición de $\operatorname{Alt}$, tenemos
\end{enumerate}
\[
\operatorname{Alt}(\theta)(v_{\tau(1)},\dots,v_{\tau(k)})=\frac{1}{k!}\sum_{\sigma\in S_k}\operatorname{sg}(\sigma)\theta(v_{\tau(\sigma(1))},\dots,v_{\tau(\sigma(k))}).
\]
\begin{enumerate}[start=6]
\item Como $v_{\tau(\sigma(r))}=v_{(\tau\circ\sigma)(r)}$, podemos escribir
\end{enumerate}
\[
\operatorname{Alt}(\theta)(v_{\tau(1)},\dots,v_{\tau(k)})=\frac{1}{k!}\sum_{\sigma\in S_k}\operatorname{sg}(\sigma)\theta(v_{(\tau\circ\sigma)(1)},\dots,v_{(\tau\circ\sigma)(k)}).
\]
\begin{enumerate}[start=7]
\item Hacemos el cambio de variable $\rho=\tau\circ\sigma$. Cuando $\sigma$ recorre $S_k$, también $\rho$ recorre $S_k$. Además, $\sigma=\tau^{-1}\circ\rho$.
\item Entonces
\end{enumerate}
\[
\operatorname{sg}(\sigma)=\operatorname{sg}(\tau^{-1}\circ\rho)=\operatorname{sg}(\tau^{-1})\operatorname{sg}(\rho)=\operatorname{sg}(\tau)\operatorname{sg}(\rho).
\]
\begin{enumerate}[start=9]
\item Sustituyendo en la suma, obtenemos
\end{enumerate}
\[
\begin{aligned}\operatorname{Alt}(\theta)(v_{\tau(1)},\dots,v_{\tau(k)})&=\frac{1}{k!}\sum_{\rho\in S_k}\operatorname{sg}(\tau)\operatorname{sg}(\rho)\theta(v_{\rho(1)},\dots,v_{\rho(k)})\\&=\operatorname{sg}(\tau)\frac{1}{k!}\sum_{\rho\in S_k}\operatorname{sg}(\rho)\theta(v_{\rho(1)},\dots,v_{\rho(k)})\\&=\operatorname{sg}(\tau)\operatorname{Alt}(\theta)(v_1,\dots,v_k).\end{aligned}
\]
\begin{enumerate}[start=10]
\item Por lo tanto, $\operatorname{Alt}(\theta)$ es alternante. Como además es multilineal, se tiene $\operatorname{Alt}(\theta)\in\Lambda^k(V^*)$.
\end{enumerate}
\begin{itemize}
\item \textbf{Linealidad}
\end{itemize}
\begin{enumerate}
\item Ahora probamos que $\operatorname{Alt}:T^k(V^*)\to\Lambda^k(V^*)$ es lineal. Sean $\theta,\eta\in T^k(V^*)$ y $a,b\in\mathbb R$.
\item Entonces
\end{enumerate}
\[
\begin{aligned}\operatorname{Alt}(a\theta+b\eta)(v_1,\dots,v_k)&=\frac{1}{k!}\sum_{\sigma\in S_k}\operatorname{sg}(\sigma)(a\theta+b\eta)(v_{\sigma(1)},\dots,v_{\sigma(k)})\\&=a\frac{1}{k!}\sum_{\sigma\in S_k}\operatorname{sg}(\sigma)\theta(v_{\sigma(1)},\dots,v_{\sigma(k)})\\&\quad+b\frac{1}{k!}\sum_{\sigma\in S_k}\operatorname{sg}(\sigma)\eta(v_{\sigma(1)},\dots,v_{\sigma(k)})\\&=a\operatorname{Alt}(\theta)(v_1,\dots,v_k)+b\operatorname{Alt}(\eta)(v_1,\dots,v_k).\end{aligned}
\]
\begin{enumerate}[start=3]
\item Por lo tanto,
\end{enumerate}
\[
\operatorname{Alt}(a\theta+b\eta)=a\operatorname{Alt}(\theta)+b\operatorname{Alt}(\eta).
\]
Así, $\operatorname{Alt}$ es lineal.

\begin{itemize}
\item $\operatorname{Alt}(\theta)=\theta$
\end{itemize}
\begin{enumerate}
\item Finalmente, supongamos que $\theta\in\Lambda^k(V^*)$. Entonces, para toda $\sigma\in S_k$, se cumple
\end{enumerate}
\[
\theta(v_{\sigma(1)},\dots,v_{\sigma(k)})=\operatorname{sg}(\sigma)\theta(v_1,\dots,v_k).
\]
\begin{enumerate}[start=2]
\item Por la definición de $\operatorname{Alt}$,
\end{enumerate}
\[
\begin{aligned}\operatorname{Alt}(\theta)(v_1,\dots,v_k)&=\frac{1}{k!}\sum_{\sigma\in S_k}\operatorname{sg}(\sigma)\theta(v_{\sigma(1)},\dots,v_{\sigma(k)})\\&=\frac{1}{k!}\sum_{\sigma\in S_k}\operatorname{sg}(\sigma)\operatorname{sg}(\sigma)\theta(v_1,\dots,v_k).\end{aligned}
\]
\begin{enumerate}[start=3]
\item Como $\operatorname{sg}(\sigma)^2=1$, resulta
\end{enumerate}
\[
\begin{aligned}\operatorname{Alt}(\theta)(v_1,\dots,v_k)&=\frac{1}{k!}\sum_{\sigma\in S_k}\theta(v_1,\dots,v_k)\\&=\frac{1}{k!}k!\theta(v_1,\dots,v_k)\\&=\theta(v_1,\dots,v_k).\end{aligned}
\]
\begin{enumerate}[start=4]
\item Como esto vale para todo $v_1,\dots,v_k\in V$, concluimos que
\end{enumerate}
\[
\operatorname{Alt}(\theta)=\theta.
\]
\bookqed
\end{bookproof}
\phantomsection\label{blk:c407fe}
\begin{bookdefinition}{Producto exterior o producto cuña}
Sean $\omega\in\Lambda^k(V^*)$, $\theta\in\Lambda^l(V^*)$. Se define el producto exterior de $\omega$ y $\theta$, denotado por $\omega\wedge\theta$, como el $(k+l)$-tensor covariante

\[
\omega\wedge\theta:=\frac{(k+l)!}{k!\,l!}\operatorname{Alt}(\omega\otimes\theta).
\]
Es decir,

\[
(\omega\wedge\theta)(v_1,\ldots,v_{k+l})= \frac{(k+l)!}{k!\,l!}\frac1{(k+l)!}\sum_{\sigma\in S_{k+l}}(\operatorname{sign}\sigma)(\omega\otimes\theta)(v_{\sigma(1)},\ldots,v_{\sigma(k+l)})
\]
\[
= \frac1{k!\,l!}\sum_{\sigma\in S_{k+l}}(\operatorname{sign}\sigma)\, \omega(v_{\sigma(1)},\ldots,v_{\sigma(k)})\theta(v_{\sigma(k+1)},\ldots,v_{\sigma(k+l)}).
\]
\end{bookdefinition}
\begin{bookproposition}{Ejercicio}
\begin{enumerate}
\item $(\omega_1+\omega_2)\wedge\theta=\omega_1\wedge\theta+\omega_2\wedge\theta$.
\item $\omega\wedge(\theta_1+\theta_2)=\omega\wedge\theta_1+\omega\wedge\theta_2$.
\item Si $s\in\mathbb R$, entonces $(s\omega)\wedge\theta=\omega\wedge(s\theta)=s(\omega\wedge\theta)$.
\item Asociativa:
\end{enumerate}
\[
\omega\wedge(\theta\wedge\xi):=(\omega\wedge\theta)\wedge\xi.
\]
\begin{enumerate}[start=5]
\item Anticonmutatividad: si $\omega\in\Lambda^k(V^*)$, $\theta\in\Lambda^l(V^*)$, entonces
\end{enumerate}
\[
\omega\wedge\theta=(-1)^{kl}\theta\wedge\omega.
\]
\begin{enumerate}[start=6]
\item Menores: si $\alpha_1,\ldots,\alpha_k\in V^*$, entonces
\end{enumerate}
\[
(\alpha_1\wedge\cdots\wedge\alpha_k)(v_1,\ldots,v_k)=\det(\alpha_i(v_j)).
\]
\begin{enumerate}[start=7]
\item Si $\theta\in\Lambda^k(V^*)$ con $k$ impar, entonces
\end{enumerate}
\[
\theta\wedge\theta=0.
\]
En particular, si $\alpha\in V^*$, entonces $\alpha\wedge\alpha=0$. 8. Leer Tu seccion 3.9. 9. Leer Tu, prop. 3.27.

\end{bookproposition}
\begin{bookproof}{Ayuda para (5)}
Sean $\omega\in\Lambda^k(V^*)$, $\theta\in\Lambda^l(V^*)$. Se define $\mu\in S_{k+l}$ por

\[
\mu=(1,\ldots,l,l+1,\ldots,l+k;\ k+1,\ldots,k+l,1,\ldots,k).
\]
Asi,

\[
\omega\wedge\theta(v_1,\ldots,v_{k+l}) =\frac1{k!\,l!}\sum_{\sigma\in S_{k+l}}(\operatorname{sign}\sigma)\omega(v_{\sigma(1)},\ldots,v_{\sigma(k)}) \theta(v_{\sigma(k+1)},\ldots,v_{\sigma(k+l)}).
\]
Resta calcular el signo de $\mu$. Considerar

\[
\rho_1=\begin{pmatrix} 1&2&\cdots&k+l-1&k+l\\ 2&3&\cdots&k+l&1 \end{pmatrix},
\]
\[
\rho_2=\begin{pmatrix} 1&2&\cdots&k+l-1&k+l\\ 3&4&\cdots&1&2 \end{pmatrix},
\]
\[
\vdots
\]
\[
\rho^k=\begin{pmatrix} 1&\cdots&l&l+1&\cdots&k+l\\ k+1&\cdots&k+l&1&\cdots&k \end{pmatrix}=\mu.
\]
Por tanto, $\operatorname{sign}\mu=(\operatorname{sign}\rho)^k$. Como $\rho$ es un ciclo,

\[
\operatorname{sign}\rho=(-1)^{k+l-1},
\]
de donde

\[
\operatorname{sign}\mu=\left((-1)^{k+l-1}\right)^k=(-1)^{kl}.
\]
\begin{itemize}
\item Ayuda para (6) Venimos a probar que $\theta$ es un conjunto L.I. si
\end{itemize}
\[
\theta=\sum_{1\le i_1<\cdots<i_k\le n}c_{i_1,\ldots,i_k}e^*_{i_1}\wedge\cdots\wedge e^*_{i_k}.
\]
Evaluamos en la $k$-upla $(e_{j_1},\ldots,e_{j_k})$, con $1\le j_1<\cdots<j_k\le n$.

Como

\[
e^*_{i_1}\wedge\cdots\wedge e^*_{i_k}(e_{j_1},\ldots,e_{j_k}) = \det(e^*_{i_p}(e_{j_q}))_{p,q=1}^k,
\]
entonces vale $1$ si $(i_1,\ldots,i_k)=(j_1,\ldots,j_k)$ y $0$ en otro caso. Luego los coeficientes son cero.

\bookqed
\end{bookproof}
\phantomsection\label{blk:8619e2}
\begin{booktheorem}{}
Sea $\{e_1,\ldots,e_n\}$ base de $V$ y $\{e^*_1,\ldots,e^*_n\}$ la respectiva base dual de $V^*$. Si $k\le n$, entonces

\[
\mathcal B=\{e^*_{i_1}\wedge\cdots\wedge e^*_{i_k}:1\le i_1<\cdots<i_k\le n\}
\]
es una base de $\Lambda^k(V^*)$ y

\[
\dim\Lambda^k(V^*)=\binom nk.
\]
Si $k>n$, entonces $\Lambda^k(V^*)=\{0\}$.

\end{booktheorem}
\begin{bookproof}{}
Veamos primero que $\mathcal B$ es un conjunto L.I. Si

\[
0=\sum_{1\le i_1<\cdots<i_k\le n}c_{i_1,\ldots,i_k}e^*_{i_1}\wedge\cdots\wedge e^*_{i_k},
\]
evaluamos en la $k$-upla $(e_{j_1},\ldots,e_{j_k})$, con $1\le j_1<\cdots<j_k\le n$.

Como

\[
e^*_{i_1}\wedge\cdots\wedge e^*_{i_k}(e_{j_1},\ldots,e_{j_k}) = \begin{cases} 1 & \text{si }(i_1,\ldots,i_k)=(j_1,\ldots,j_k),\\ 0 & \text{en otro caso}, \end{cases}
\]
resulta $c_{i_1,\ldots,i_k}=0$ para todo $1\le i_1<\cdots<i_k\le n$.

Veamos que genera. Sea $\theta\in\Lambda^k(V^*)$ y entonces, como $\theta\in T^k(V^*)$,

\[
\theta=\sum a_{i_1,\ldots,i_k}e^*_{i_1}\otimes\cdots\otimes e^*_{i_k}.
\]
Entonces, como $\theta\in\Lambda^k(V^*)$,

\[
\theta=\operatorname{Alt}\theta= \sum a_{i_1,\ldots,i_k}\operatorname{Alt}(e^*_{i_1}\otimes\cdots\otimes e^*_{i_k}).
\]
Luego se reordena cada termino, y si hay al menos dos indices iguales el termino es cero; si no,

\[
\operatorname{Alt}(e^*_{i_1}\otimes\cdots\otimes e^*_{i_k})=\text{multiplo de }e^*_{i_{\sigma(1)}}\wedge\cdots\wedge e^*_{i_{\sigma(k)}}
\]
para indices ordenados. Por lo que genera.

\bookqed
\end{bookproof}
\begin{bookdefinition}{Algebra exterior}
Se define $\Lambda^0(V^*)=\mathbb R$. Se define

\[
\Lambda(V^*):=\Lambda^0(V^*)\oplus\Lambda^1(V^*)\oplus\cdots\oplus\Lambda^n(V^*)
\]
y se extiende el producto cuña $\wedge$ a todo $\Lambda(V^*)$ de manera bilineal. Es decir, si

\[
\alpha=\alpha_0+\alpha_1+\cdots+\alpha_n,\qquad \alpha_i\in\Lambda^i(V^*)
\]
y

\[
\beta=\beta_0+\beta_1+\cdots+\beta_n,\qquad \beta_j\in\Lambda^j(V^*),
\]
entonces se define

\[
\alpha\wedge\beta:=\sum_{i=0}^n\sum_{j=0}^n\alpha_i\wedge\beta_j,
\]
donde

\[
t\in\Lambda^0(V^*),\quad \theta\in\Lambda^k(V^*)\implies t\wedge\theta=t\theta.
\]
Por las propiedades anteriores, $(\Lambda(V^*),\wedge)$ es una algebra asociativa, no conmutativa y con unidad $1\in\Lambda^0(V^*)$, que tiene dimension

\[
\binom n0+\binom n1+\cdots+\binom nn=2^n,
\]
y se llama algebra exterior.

\end{bookdefinition}
\phantomsection\label{blk:2c0c65}
